\clearpage
\section{Corrected rotation algebra and support estimates}\label{app:rotation}
This appendix records the algebra, normalization and domain extension used in Section~\ref{sec:rotation}. The rotation anticommutator and kernel-invariance theorem are due to the cited prior work; the quantitative use of a bounded internal radius is the estimate needed here.

\subsection{Two charge conventions and two summation scopes}
Retain the kinetic-one Hamiltonian $h$. Let $\widetilde q$ denote the convention
\begin{equation}
 \{\widetilde q_\alpha,\widetilde q_\beta\}=\delta_{\alpha\beta}h,
 \qquad\sum_{\alpha=1}^{16}\norm{\widetilde q_\alpha g}^2=8\hh[g]
\end{equation}
on the physical core. The manuscript's baseline charges are $q=\sqrt2\widetilde q$ and have stack coefficient $16$. For a fixed $j=(u,v)$ the primitive of \cite[Lemma 2(c)]{HHI} is
\begin{equation}
 \widetilde a_{j,\alpha}=\frac1{16}(X_u\Gamma_v-X_v\Gamma_u)_{\alpha\epsilon}\theta_\epsilon
 +\frac1{56}\sum_{w\notin\{u,v\}}X_w(\Gamma_w\Gamma_{uv})_{\alpha\epsilon}\theta_\epsilon.
 \label{eq:expandedprimitive}
\end{equation}
The color index is contracted in each summand. This follows from $(5b+2c)/7$ because $\Gamma_u\Gamma_{uv}=\Gamma_v$ and $\Gamma_v\Gamma_{uv}=-\Gamma_u$.

All coefficient matrices are real. The first two are symmetric grade-one matrices, the others antisymmetric grade-three matrices. They are mutually orthogonal in the Frobenius inner product and each has squared Frobenius norm $16$. Each $\widetilde a_{j,\alpha}$ is a real linear combination of Majoranas, so its square is scalar by the CAR. Consequently
\begin{align}
 \sum_\alpha\widetilde a_{j,\alpha}^2
 &=\frac12\left[\frac{16}{16^2}(|X_u|^2+|X_v|^2)
          +\frac{16}{56^2}\sum_{w\notin\{u,v\}}|X_w|^2\right]I\\
 &=\left[\frac{|X_u|^2+|X_v|^2}{32}
          +\frac{\sum_{w\notin\{u,v\}}|X_w|^2}{392}\right]I.
\end{align}
For several blocks, concatenate color coordinates; no color-dimension factor occurs. A fixed coordinate belongs to eight of the $36$ pairs and is outside the other $28$, giving
\begin{equation}
 \sum_{u<v}\sum_\alpha\widetilde a_{(u,v),\alpha}^2
 =\left(\frac8{32}+\frac{28}{392}\right)\alpha^2I
 =\frac9{28}\alpha^2I.\label{eq:allsum}
\end{equation}
Thus $1/32$ is the bound for one generator, while $9/28$ is the exact coefficient after summing all generators. They are not alternative estimates for the same sum.

Rescale reciprocally:
\begin{equation}
 q_\alpha=\sqrt2\widetilde q_\alpha,\qquad
 a_{j,\alpha}=\widetilde a_{j,\alpha}/\sqrt2,\qquad
 J_j=\sum_\alpha\{q_\alpha,a_{j,\alpha}\}.
\end{equation}
The fixed-generator bound becomes $\alpha^2/64$, and \eqref{eq:allsum} becomes $9\alpha^2/56$. The products $8(1/32)$ and $16(1/64)$ agree. Rescaling the charge alone would change the generator identity and is not allowed.

\subsection{Direct Clifford verification of the identity}
Write $\widetilde a_\alpha=X_w(M_w)_{\alpha\epsilon}\theta_\epsilon$ for a fixed pair $(u,v)$. In the real gamma convention the coefficient matrices obey
\begin{align}
 \Tr(\Gamma_t^TM_w)&=\delta_{wu}\delta_{tv}-\delta_{wv}\delta_{tu},\\
 \sum_w\Gamma_w^TM_w&=\tfrac14\Gamma_{uv},\\
 \Tr(M_w^T\Gamma_{st})&=0.
\end{align}
The first identity yields the orbital rotation, the second the spin contribution $-i\theta\Gamma_{uv}\theta/4$, and the last removes the interaction anticommutator. These are exact finite Clifford identities. The included integer-matrix program checks them for all $36$ pairs and all indicated indices. It uses a self-contained real symmetric Spin(9) representation constructed from the octonion multiplication table. This test verifies the finite identities only; integration by parts and closed-domain statements require the arguments in the text.

\subsection{A stronger Casimir variant}
The weaker coefficient $1/4$ in Lemma~\ref{lem:ball} is sufficient for the manuscript. For comparison, the all-generator identity proves
\begin{equation}
 H_I[g]\ge\frac7{9R^2}\norm{(1-P_I)g}^2
 \qquad\text{when }\supp g\subset\{\alpha\le R\}.\label{eq:casimirball}
\end{equation}
Use the graded sum of block supercharges, or equivalently a blockwise charge index, with stack coefficient $8$. Its squares add to the complete $H_I$. Let $C_I=\sum_j(J_j^I)^2$ and initially take a physical smooth vector in a Casimir eigenspace $C_Ig=\lambda g$, $\lambda>0$. Put $E=H_I[g]$ and $K=9/28$. Rotational invariance of the complete form gives
\begin{equation}
 \sum_j\norm{J_j^Ig}^2=\lambda\norm g^2,
 \qquad\sum_jH_I[J_j^Ig]=\lambda E.
\end{equation}
Individual supercharges need not commute with rotations; only their averaged square is used. Pair the rotation identity with $J_j^Ig$, sum over $j$, and integrate by parts on the core. The first resulting term is bounded by
$\sqrt{8\lambda E}\sqrt{KR^2\norm g^2}$. For the other, pointwise Cauchy--Schwarz and the fact that each primitive square is scalar give
\begin{equation}
 \sum_\alpha\left\lVert\sum_j\widetilde a_{j,\alpha}(X)(J_j^Ig)(X)\right\rVert^2
 \le K\alpha(X)^2\sum_j\lVert(J_j^Ig)(X)\rVert^2.
\end{equation}
That term has the same bound. Therefore
\begin{equation}
 \lambda\norm g^2\le2\sqrt{8K}\,R\sqrt{\lambda E}\norm g,
 \qquad E\ge\frac{7\lambda}{72R^2}\norm g^2.
\end{equation}
In the generator convention used here, a dominant $B_4$ weight has Casimir
\begin{equation}
 \lambda_C=\sum_{i=1}^4\lambda_i(\lambda_i+9-2i),\qquad
 \lambda_1\ge\lambda_2\ge\lambda_3\ge\lambda_4\ge0,
\end{equation}
with all coordinates integral or all half-integral. The smallest positive integer case $(1,0,0,0)$ has Casimir $8$, and the smallest half-integral case $(1/2,1/2,1/2,1/2)$ has Casimir $9$. Summing positive Casimir sectors proves \eqref{eq:casimirball}.

The rotation-invariant ball preserves all compact-group projections. At finite Casimir cutoff the generators are bounded angular operators, and the projected smooth physical core remains radially compact. Invariant nonnegative forms split orthogonally over the Casimir decomposition. Finite cutoff, approximation in the form norm, and a limit therefore extend the result to the closed form domain. For support in a closed ball use a slightly larger open radius and decrease it to $R$. Spectator Hilbert spaces pass through the same argument by finite sums and closure. Neither $\widetilde a g$ nor $Jg$ is asserted to lie in the Hamiltonian operator domain for a general form vector.

\subsection{Kernel invariance without an assumed positive moment}
For completeness, the domain mechanism behind \cite[Theorem 1(a)]{HHI} is as follows. The physical kernel is invariant under the compact rotation group. Decompose it into isotypes so that $J_j$ is bounded on each finite-dimensional representation factor. Choose a smooth radial cutoff $\chi_R$ that is one in a ball and whose derivative is supported where $R<\rho<3R$. In the generator identity, test with two kernel vectors of finite isotypic type and move the charge through $\chi_R$. The surviving shell terms have bounds of the form
\begin{equation}
 C\norm{\phi}_{\{R<\rho<3R\}}
  \norm{\Psi}_{\{R<\rho<3R\}},
\end{equation}
since $[\widetilde q,\chi_R]=O(R^{-1})$ and $\widetilde a=O(R)$ on that shell. The compatible maximal/self-adjoint realization and physical-core approximation are the ones in the cited theorem. The bound tends to zero by $L^2$ integrability alone. Taking $\phi=J_j\Psi$ within each isotype gives $J_j\Psi=0$. Completing the compact-group decomposition proves invariance of the full kernel. This is an explanation of the cited prior theorem, not a new existence or domain theorem.
